
% =====================================================================
\section{Numerical verification}
\label{sec:numerical}
% =====================================================================

The theorems are proved from the definitions and require no computation. It is
nonetheless instructive to exhibit them on concrete instances of the hypothesis,
both as a check against error and to display the quantitative shape of each
result. This section reports an exhaustive numerical witness: every theorem-group
is tested on a large family of explicitly constructed bounded resolvable spaces,
and the measured quantities are plotted against the bounds the theorems predict.

\subsection{Method}

A bounded resolvable space is realised as a finite metric--measure space: a finite
integer-lattice box, a counting measure with a fixed atom weight, and a partition
into axis-aligned cells of side $\ge\reso$. By construction such an instance
satisfies \Cref{ax:brs}---it is bounded with finite measure, every cell has
diameter at least the resolution, and finer boxes refine it---and an embedded
agent on it can name only unions of cells, exactly the operational restriction of
\Cref{def:realisable}. On each instance the realisable regions, their separators,
boundary thickness, scale, log-measure, and the committed record are computed
directly from the definitions, and each theorem's inequality or equivalence is
checked.

A randomised sweep drew instances of varying dimension, extent, cell size, and
atom weight, and ran all seven theorem-group checks on each. Across
$500$ instances the suite performed $96{,}024$ individual checks, all of which
passed: no realisable separator was sharp; every boundary thickness met
$\thick\ge\mu_{\min}$; the three floor derivations agreed on every instance;
distinguishability matched $\sigma>\thick_A$ in every region tested; the
self-referential value admitted no consistent assignment; and the committed
record was monotone on every traversal. The figures below are drawn from these
same computations.

\subsection{Boundary thickness and the floor}

\Cref{fig:panel1} verifies the Boundary--Thickness Theorem. Every measured
thickness lies on or above the line $\thick_\Part(A)=\mu_{\min}$
(\Cref{thm:thickness}); the distribution of the ratio $\thick/\mu_{\min}$ is
supported entirely at and above $1$, so no realisable separator is sharp; and the
minimum thickness over each instance, plotted against cell count and
$\mu_{\min}$, stays strictly positive throughout.

\begin{figure}[h]
\centering
\includegraphics[width=\textwidth]{figures/panel_1.png}
\caption{Boundary--Thickness Theorem (\Cref{thm:thickness}). Left to right:
measured thickness against $\mu_{\min}$ with the equality line; histogram of
$\thick/\mu_{\min}$, entirely $\ge1$; thickness against region scale; and the
per-instance minimum thickness over cell count and $\mu_{\min}$.}
\label{fig:panel1}
\end{figure}

\Cref{fig:panel2} verifies that the three derivations bound one constant
(\Cref{thm:floor-agree}): the geometric floor agrees with the representational
and cost floors instance by instance (points on the diagonal), and the cost
$g(t)=-\log t$ diverges as $t\to0^+$, so no finite-cost process reaches the sharp
cut.

\begin{figure}[h]
\centering
\includegraphics[width=\textwidth]{figures/panel_2.png}
\caption{Agreement of the three floor derivations (\Cref{thm:floor-agree}).
Geometric versus representational floor; geometric versus cost floor; the
divergence of $g(t)=-\log t$; and the joint surface of the three floors.}
\label{fig:panel2}
\end{figure}

\subsection{Detectability and the sorites}

\Cref{fig:panel3} verifies the Detectability Theorem. A region carries positive
interior content exactly on the side $\sigma>\thick_A$ of the diagonal, and the
interior content is positive precisely when the margin $\sigma-\thick_A$ is
positive (\Cref{thm:detectability}); the distribution of $\tau=\thick/\sigma$
separates the engulfed regime $\tau\ge1$ from the distinguishable regime
$\tau<1$.

\begin{figure}[h]
\centering
\includegraphics[width=\textwidth]{figures/panel_3.png}
\caption{Detectability (\Cref{thm:detectability}). Scale against $\thick_A$ with
distinguishable and engulfed regions separated by the diagonal; interior content
against the margin $\sigma-\thick_A$; the distribution of $\tau=\thick/\sigma$
about the threshold $1$; and the joint surface of scale, $\thick_A$, and
interior content.}
\label{fig:panel3}
\end{figure}

\Cref{fig:panel4} verifies the sorites corollary. As a region grows by sub-floor
grains, $\tau$ descends smoothly through the band around $1$; the per-step change
in $\tau$ stays bounded and decays once the aggregate clears its boundary, so no
single grain carries $\tau$ across the threshold (\Cref{cor:sorites}). Three
independent accretion orders all descend through the same band.

\begin{figure}[h]
\centering
\includegraphics[width=\textwidth]{figures/panel_4.png}
\caption{The sorites dissolved (\Cref{cor:sorites}). The relative thickness
$\tau$ against grains added and against scale, both crossing $1$ across a band;
the bounded per-step change $|\Delta\tau|$; and three accretion orders descending
through the band.}
\label{fig:panel4}
\end{figure}

\subsection{Locating versus visiting}

\Cref{fig:panel5} verifies the Locating--Visiting Asymmetry. The locating cost
rises in steps of $\thick$ as $\thick\lceil\log_2 K\rceil$ while the visiting
cost is identically zero (\Cref{thm:asymmetry}, \Cref{thm:locating-cost}); the
cost is linear in $\log_2(1/p)$; and the visiting frequency equals $p(W)$, so
rarity enters only through frequency.

\begin{figure}[h]
\centering
\includegraphics[width=\textwidth]{figures/panel_5.png}
\caption{Locating versus visiting (\Cref{thm:asymmetry},
\Cref{thm:locating-cost}). Locating cost and (zero) visiting cost against the
number of alternatives $K$; locating cost against $\log_2(1/p)$; visiting
frequency equal to $p(W)$; and the locating-cost surface over $\log_2 K$ and the
floor $\thick$.}
\label{fig:panel5}
\end{figure}

\subsection{Non-return and cessation}

\Cref{fig:panel6} verifies the closing results. The committed record is monotone
and an attempted undo still increments it (\Cref{thm:nonreturn}); the cost
$g(t)$ diverges while the return is bounded, so the marginal return per cost falls
to zero and a positive halting thickness $\thick_\ast$ exists for every positive
threshold (\Cref{thm:cessation}).

\begin{figure}[h]
\centering
\includegraphics[width=\textwidth]{figures/panel_6.png}
\caption{Non-return and cessation (\Cref{thm:nonreturn}, \Cref{thm:cessation}).
The monotone committed record with the undo step marked; the diverging cost
$g(t)$; the marginal return per cost with several halting thresholds
$\thick_\ast$; and the marginal-return surface over thickness and threshold.}
\label{fig:panel6}
\end{figure}

\begin{remark}[Status of the numerical witness]
\label{rem:numerical-status}
The figures confirm the theorems on concrete instances and display their
quantitative form; they are evidence of correctness and of the magnitudes
involved, not a substitute for the proofs, which hold for every space satisfying
\Cref{ax:brs}. That every one of the $96{,}024$ checks across $500$ independently
constructed instances agreed with the predicted bounds is consistent with the
claim that the floor is structural rather than an artefact of any one
construction.
\end{remark}
